
\documentclass[11pt]{article}
\usepackage[a4paper,margin=28mm]{geometry}
\usepackage{amsmath,amssymb,amsthm}
\usepackage{enumitem}
\usepackage[hidelinks]{hyperref}

\newtheorem{definition}{Definition}
\newtheorem{conjecture}{Conjecture}
\newtheorem{question}{Research Question}
\newtheorem{theorem}{Theorem}

\title{Hard Problem: The Generalized Star-Height Conjecture}
\author{Yoav Danieli}
\date{}

\begin{document}
\maketitle

\section{Ordinary regular languages}

Fix a finite alphabet $\Sigma$.  Write $\Sigma^*$ for the set of finite words
over $\Sigma$, including the empty word $\varepsilon$.

\begin{definition}[Generalized regular expression]
The set of generalized regular expressions over $\Sigma$ is generated by
\[
 E ::= \emptyset\mid\varepsilon\mid a
 \mid E\cup E\mid E\cdot E\mid E^*\mid \neg E,
 \qquad a\in\Sigma.
\]
The semantics is:
\[
\begin{aligned}
L(\emptyset)&=\emptyset,&
L(\varepsilon)&=\{\varepsilon\},&
L(a)&=\{a\},\\
L(E\cup F)&=L(E)\cup L(F),&
L(E\cdot F)&=\{uv:u\in L(E),v\in L(F)\},\\
L(E^*)&=\bigcup_{n\geq0}L(E)^n,&
L(\neg E)&=\Sigma^*\setminus L(E).
\end{aligned}
\]
Thus complement is always taken relative to the fixed universe $\Sigma^*$.
Intersection and set difference are definable from union and complement.
\end{definition}

\begin{definition}[Star height of an expression]
The generalized star height $h(E)$ is the maximum number of Kleene-star
operators on a path from the root of the syntax tree of $E$ to a leaf:
\[
\begin{aligned}
h(\emptyset)=h(\varepsilon)=h(a)&=0,\\
h(\neg E)&=h(E),\\
h(E\cup F)=h(E\cdot F)&=\max\{h(E),h(F)\},\\
h(E^*)&=1+h(E).
\end{aligned}
\]
\end{definition}

\begin{definition}[Generalized star height of a language]
For a regular language $L\subseteq\Sigma^*$, define
\[
 h_{\mathrm g}(L)
 =
 \min\{h(E):E\text{ is a generalized regular expression and }L(E)=L\}.
\]
The minimum exists because every regular language has an ordinary regular
expression, which is also a generalized regular expression.
\end{definition}

A language has generalized star height $0$ precisely when it can be defined
without Kleene star; such languages are called \emph{star-free}.  Star-free
regular languages admit the classical algebraic characterization by
aperiodic syntactic monoids.

\section{Exact open question}

\begin{question}[Generalized star-height problem]
Does there exist a regular language $L$ such that
\[
  h_{\mathrm g}(L)\geq2?
\]
\end{question}

This question is equivalent to the following conjecture.

\begin{conjecture}[Generalized star-height conjecture]
For every finite alphabet $\Sigma$ and every regular language
$L\subseteq\Sigma^*$,
\[
  h_{\mathrm g}(L)\leq1.
\]
\end{conjecture}


\begin{thebibliography}{9}

\bibitem{Schuetzenberger}
M.-P.~Sch\"utzenberger,
\emph{On Finite Monoids Having Only Trivial Subgroups},
Information and Control 8 (1965), 190--194.

\bibitem{Pin}
J.-E.~Pin,
\emph{Some Results on the Generalized Star-Height Problem},
Information and Computation 101 (1992), 219--250.

\bibitem{Sakarovitch}
J.~Sakarovitch,
\emph{Elements of Automata Theory},
Cambridge University Press, 2009.

\end{thebibliography}

\end{document}
